% Part: lambda-calculus
% Chapter: syntax
% Section: eta

\documentclass[../../../include/open-logic-section]{subfiles}

\begin{document}

\olfileid[pt]{lam}{syn}{eta}
\olsection{Conversão~$\eta$}

Há outra relação sobre termos~$\lambda$. Em \olref[fv]{sec}, usamos o exemplo
$\lambd[x][(fx)]$, que recebe um argumento e aplica $f$ a ele. Em outras
palavras, é a mesma função que~$f$: tanto $\lambd[x][(fx)]N$ como $fN$ se
reduzem a~$fN$. Usamos a redução~$\eta$ (e a extensão~$\eta$) para capturar
essa ideia.

\begin{defn}[Contração~$\eta$, $\eredone$]
  \ollabel{defn:beredone}
  A \emph{contração~$\eta$} ($\eredone$) é a menor relação compatível sobre
  termos que satisfaz a seguinte condição:
  \[
    \lambd[x][M x] \eredone M \text{ desde que } x \notin FV(M)
  \]
\end{defn}

\begin{defn}[Redução~$\beta\eta$, $\bered$] \ollabel{defn:bered}
  A \emph{redução~$\beta\eta$} ($\bered$) é a menor relação reflexiva e
  transitiva sobre termos que contém $\bredone$ e $\eredone$; isto é, as
  regras de reflexividade e transitividade mais as duas regras a seguir:
  \begin{enumerate}
  \item Se $M \bredone N$, então $M \bered N$. \ollabel{defn:bered3}
  \item Se $M \eredone N$, então $M \bered N$. \ollabel{defn:bered4}
  \end{enumerate}
    
\end{defn}

\begin{defn}
  Ampliamos a relação de equivalência $\equal$ com a regra de conversão~$\eta$:
  \[
  \lambd[x][f x] \equal f
  \]
  e denotamos a relação ampliada por $\equal[\eta]$.
\end{defn}

A equivalência~$\eta$ é importante porque está relacionada à extensionalidade
dos termos lambda:

\begin{defn}[Extensionalidade]
  Ampliamos a relação de equivalência $\equal$ com a regra~(\ext):
  \begin{center}
  Se $Mx \equal Nx$, então $M \equal N$, desde que $x \notin FV(MN)$.
  \end{center}
  e denotamos a relação ampliada por $\equal[\ext]$.
\end{defn}

Em linhas gerais, a regra afirma que dois termos, considerados como funções,
devem ser considerados iguais se se comportam da mesma maneira para o mesmo
argumento.

Mostraremos agora que a regra~$\eta$ fornece exatamente a extensionalidade, e
nada mais.

\begin{thm}
  $M \equal[\ext] N$ se e somente se $M \equal[\eta] N$.
\end{thm}

\begin{proof}
  Primeiro mostramos que $\equal[\eta]$ é fechada sob a regra de
  extensionalidade. Isto é, a regra $ext$ não acrescenta nada a
  $\equal[\eta]$. Temos, então, que $\equal[\eta]$ contém $\equal[\ext]$ e,
  se $M \equal[\ext] N$, então $M \equal[\eta] N$.

  Para mostrar que $\equal[\eta]$ é fechada sob~\ext, observe que, para todo
  $M \equal N$ derivado pela regra~\ext{}, temos $Mx \equal[\eta] Nx$ como
  premissa. Então temos $\lambd[x][Mx] \equal[\eta] \lambd[x][Nx]$ por uma
  regra de~$\equal$; aplicando~$\eta$ aos dois
  lados, obtemos $M \equal[\eta] N$.

  De modo análogo, mostramos que a regra~$\eta$ está contida em
  $\equal[ext]$. Para quaisquer $\lambd[x][Mx]$ e $M$ com $x \notin FV(M)$,
  temos $(\lambd[x][Mx])x \equal[\ext] Mx$, o que nos dá
  $\lambd[x][Mx] \equal[\ext] M$ pela regra~\ext{}.
\end{proof}

\end{document}
